\documentclass[a4paper,8pt]{article}

\usepackage{parskip} 
\usepackage{hologo}
\usepackage{fontspec}

% Other packages for formatting
\usepackage{graphicx}
\usepackage[usenames,dvipsnames]{xcolor}

\usepackage[left=1cm, right=1cm, top=1cm, bottom=1cm, scale=1]{geometry}
\definecolor{371e77}{HTML}{371E77}

% Tabularx environment
\usepackage{tabularx}

% For lists
\usepackage{enumitem}

% Centered version of 'X' column
\newcolumntype{C}{>{\centering\arraybackslash}X} 

% Custom section formatting
\usepackage{titlesec}
\titleformat{\section}{\Large\scshape\raggedright}{}{0em}{}[\titlerule]
\titlespacing{\section}{0pt}{2pt}{2pt}

% Hyperlinks
\usepackage[unicode]{hyperref}
\usepackage[normalem]{ulem}
\hypersetup{
    colorlinks=true,
    linkcolor=., 
    urlcolor=.,
    citecolor=black
}
\newcommand{\ulhref}[2]{\href{#1}{\uline{#2}}}

% Social icons
\usepackage{fontawesome5}

\setmainfont{Arial} 

\begin{document}
\pagestyle{empty} % no page numbers

% ==================== Header ====================
\begin{tabularx}{\linewidth}{@{} C @{}}
\color[HTML]{1C033C}\Huge{\textbf{SANJANA RANI}} \\
\color[HTML]{371E77}\large{Graduate Aerospace Engineer | Aerodynamics • Acoustics • Experimental Testing • Data Analysis} \\
\textcolor[HTML]{371E77}{
  \raisebox{-0.2\height}{\includegraphics[height=12pt]{mail.png}} \, sanjanakamboj98@gmail.com \, | \,
  \raisebox{-0.2\height}{\includegraphics[height=12pt]{phone.png}} \, +49(0)15251347327 \, | \,
    \raisebox{-0.2\height}{\includegraphics[height=12pt]{linkedin.png}} \, 
\ulhref{https://linkedin.com/in/sanjana-rani}{linkedin.com/in/sanjana-rani}
}
\end{tabularx}

\vspace{0.5cm}

% ==================== Date & Recipient ====================
\noindent
\textbf{Date:} \today \\

German Aerospace Center (DLR)\\
Recruitment Team\\
Cochstedt, Germany\\

% ==================== Subject ====================
\noindent
\textbf{Subject: Application for Research Associate --- Aircraft Noise from Advanced Air Mobility (Ref.\ No.\ 4364)}\\

Dear Hiring Team,\\

I am writing to apply for the Research Associate position in Aircraft Noise from Advanced Air Mobility (Ref. No. 4364) at the Institute of Aerodynamics and Flow Technology in Cochstedt. I hold a Joint Master’s degree in Transfer of Fluid Materials in Aeronautical \& Space Applications (TFM-ASA) and have gained hands-on experience in aeroacoustic measurements, experimental aerodynamics, and Python-based signal processing through my work at DLR and university laboratories. Having contributed to acoustic sensor validation within DLR, I am motivated to extend this experience towards aircraft noise modelling, simulation benchmarking, and advanced air mobility research.\\

During my Graduate Research Assistantship at the DLR Institute of Electrified Aero Engines, I contributed to the development and experimental validation of acoustic and vibroacoustic sensor systems for EMC/EMI testing under the LU(FT)\(^2\) 2030 project, including building a prototype EMI testing rig at the Braunschweig facility. I monitored time-based signals using DEWESoft and designed 
Python-based signal processing pipelines --- FFT, coherence analysis, and spectral subtraction --- to evaluate sensor sensitivity and noise susceptibility across 2 microphone types and 4 spatial configurations over a 0--50\,kHz frequency range. I also performed cross-modal comparisons of accelerometer and microphone magnitude spectra to characterise near-field and far-field noise behaviour and quantify measurement uncertainty. This work gave me end-to-end experience with acoustic measurement campaigns --- from rig design and data acquisition through to signal processing and uncertainty analysis.\\

Prior to DLR, as a Student Assistant at the Chair of Aerodynamics and Fluid Mechanics at BTU Cottbus, I supported the TEKUS/DEPIK zero-gravity convection experiment, contributing to mechanical assembly, experimental integration, flow-quality validation, and PIV-based flow analysis. Alongside my professional experience, my Joint MSc provided strong foundations in experimental aerodynamics 
and flow measurement techniques, complemented by project work spanning stereo-PIV flow analysis, reduced-order aerodynamic modelling in Python, and hands-on acoustic measurement --- including an SPL survey of a running compressor to characterise near-field noise distribution at varying heights. These experiences across experimental and computational workflows have strengthened my ability to independently analyse complex data while contributing effectively within collaborative research teams.\\

This position represents an opportunity to build directly on my experimental and signal processing background in a research environment I already know and respect. I am drawn to DLR's work on Advanced Air Mobility noise not only because it connects to my previous acoustic measurement experience, but because it sits at the intersection of a genuinely open scientific problem and a rapidly evolving field. I am eager to contribute to flight test campaigns and simulation benchmarking at Cochstedt, and to develop deeper expertise in noise source modelling and immission prediction methods. Longer term, I am motivated to grow as an independent researcher --- publishing results, collaborating with project partners, and contributing to DLR's position at the forefront of aeroacoustics research.\\

I am aware that the role requires strong German language skills. My current level is A1, and I am actively improving, with the goal of reaching B1 by the end of 2026. I am confident in my ability to work effectively in English in international research environments and committed to further developing my German alongside my work. I would welcome the opportunity to discuss my application further and am available for an interview at your convenience. Thank you for considering my application.\\

Sincerely,\\
\textbf{Sanjana Rani}

\end{document}
